%%--------------------------------------------------------------------------
%%  SECTION 2: STRUCTURE-PRESERVING AI FOR KINETIC EQUATIONS
%%--------------------------------------------------------------------------
\section{Structure-Preserving AI for Multiscale Kinetic Equations}

\noindent\textit{Why it matters.}\enspace
Kinetic equations are among the most computationally expensive models in all
of science.  A single RTE simulation in nuclear reactor transport involves
$10^8$--$10^{12}$ phase-space unknowns, with run times on classical
supercomputers measured in hours.  Fusion plasma simulations face a starker
challenge: real-time feedback control of a tokamak plasma requires millions
of Vlasov--Maxwell evaluations per second, a throughput that is five to six
orders of magnitude beyond what current particle-in-cell (PIC) codes can
deliver.  Equally important is the \emph{multiscale} structure of these
equations: in the limit $\varepsilon \to 0$ (small Knudsen number), the
kinetic model degenerates to a macroscopic diffusion or fluid system, but
classical numerical discretizations become stiff and impose a time-step
restriction $\Delta t = O(\varepsilon)$, adding an entire order of
computational difficulty to an already intractable problem.

\subsection{Asymptotic-Preserving Neural Networks}

The central contribution of my group in this area is the development of
\emph{Asymptotic-Preserving Neural Networks}
(APNNs)~\cite{APNN2023,APNN2024}: a fundamentally new class of
physics-informed neural network architectures that provably inherit the
asymptotic-preserving (AP) property of the best classical schemes.  The
essential novelty is architectural.  We encode the micro--macro decomposition
$f = \rho M + \varepsilon g$ (where $M$ is the local Maxwellian equilibrium
and $g$ is the microscopic deviation) directly into the network's structural
inductive bias, rather than including it as a soft loss-function term that
can be violated during optimization.  The ``macro'' branch of the network is
constrained to satisfy the correct limiting diffusion equation as
$\varepsilon \to 0$, and the ``micro'' branch is constrained to vanish in the
same limit.  This is, to our knowledge, the first class of neural network
methods for kinetic equations with a rigorous, verifiable AP guarantee
independent of training quality or random seed.

The framework has been applied across the canonical multiscale kinetic
systems:
\begin{itemize}
  \item \emph{Linear transport equations}~\cite{APNN2023,APNN2024}: uniform
    accuracy from the kinetic regime to the diffusion limit without mesh
    refinement, on problems where standard PINNs fail completely for
    $\varepsilon < 10^{-3}$.

  \item \emph{Vlasov--Poisson--Fokker--Planck system in the high-field
    regime}~\cite{APNN-VPFP2024}: AP behavior in the presence of long-range
    electromagnetic self-interactions, extending the framework to
    nonlinearly coupled systems.

  \item \emph{Gray radiative transfer via even--odd decomposition}
    (AP-CON)~\cite{APCON2024}: operator-learning variants using convolutional
    DeepONets that generalize across $10^4$ distinct optical parameter
    configurations without retraining, achieving accuracy competitive with
    reference $S_N$ solvers while running $10^3\times$ faster at inference.
\end{itemize}

\noindent Reference implementations are publicly available at
\href{https://github.com/mazhengcn/apnns}{\texttt{github.com/mazhengcn/apnns}}
and
\href{https://github.com/mazhengcn/apcons}{\texttt{github.com/mazhengcn/apcons}},
providing training scripts, pre-computed datasets, and documented APIs for
reproducing all published results and adapting the methods to new equations.

\subsection{DeepRTE: A Foundation Model for Radiative Transfer}

A more ambitious development is \textbf{DeepRTE}~\cite{DeepRTE2026}, the
first pre-trained, attention-based \emph{neural operator foundation model}
for the steady-state RTE.  Inspired by the large language model paradigm, the
idea is to train a single model once on a large, systematically generated
dataset of RTE solutions---covering the full range of physically relevant
scattering anisotropy parameters $g \in [0.1, 0.8]$---and then deploy it for
instant inference on any new optical geometry without retraining.  DeepRTE
learns the solution operator $\mathcal{A}: (I^{-};\, \mu_t, \mu_s, p)
\mapsto I$, where $I^{-}$ is the boundary inflow, $\mu_t, \mu_s$ are
attenuation and scattering coefficients, and $p$ is the scattering phase
function, mapping this entire configuration to the full radiation field.

The engineering realization is as important as the scientific contribution.
DeepRTE is implemented in JAX/Flax for hardware-accelerated training and
inference, containerized with Docker for reproducible deployment, and
distributed as a complete package at
\href{https://github.com/mazhengcn/deeprte}{\texttt{github.com/mazhengcn/deeprte}},
with pre-trained model checkpoints and training datasets hosted on Hugging
Face at
\href{https://huggingface.co/mazhengcn/deeprte}{\texttt{mazhengcn/deeprte}}.
The package provides a full training and inference pipeline, enabling
researchers at other institutions to fine-tune the model on domain-specific
geometries---nuclear shielding calculations, photon dosimetry, atmospheric
radiative transfer---without access to our cluster.  This represents a
qualitative shift: rather than a solver for individual instances, DeepRTE is
a \emph{reusable computational asset} that amortizes massive training cost
across the entire scientific community.

\subsection{Asymptotic-Preserving Random Feature Methods}

Beyond neural networks, in collaboration with Jingrun~Chen (USTC), we
developed an \emph{asymptotic-preserving random feature method}
(AP-RFM)~\cite{AP-RFM2025} for multiscale radiative transfer.  Random
feature methods expand the solution in a random basis
$\sigma(\omega_j \cdot x + b_j)$ with fixed random weights, offering
mesh-free approximation with theoretically controllable error at
$O(N)$ complexity.  By encoding the micro--macro decomposition into both
the basis construction and the least-squares fitting, AP-RFM achieves uniform
accuracy across all Knudsen numbers at a cost substantially lower than neural
network training.  This result demonstrates that the AP design principle is
a general mathematical framework---not a property specific to any particular
machine learning architecture.
